\textbf{1. Kiến thức cần nhớ}
\begin{itemize}
\item Để giải hệ phương trình bậc nhất hai ẩn, ta thường dùng hai quy tắc biến đổi tương đương:
\begin{enumerate}
\item \textbf{Phương pháp thế}: Rút một ẩn từ một phương trình rồi thế vào phương trình còn lại để thu được phương trình bậc nhất một ẩn.
\item \textbf{Phương pháp cộng đại số}: Nhân hai vế của mỗi phương trình với số thích hợp sao cho hệ số của một ẩn bằng nhau (hoặc đối nhau), rồi trừ (hoặc cộng) từng vế hai phương trình để triệt tiêu ẩn đó.
\end{enumerate}
\item \textbf{Hệ phương trình quy về bậc nhất hai ẩn}: Sử dụng phương pháp đặt ẩn phụ đưa hệ phức tạp về hệ bậc nhất hai ẩn đơn giản.
\end{itemize}

\textbf{2. Phương pháp giải các dạng}

%% DANG 9D1?2-1
\textbf{Dạng 1. Giải hệ bằng phương pháp thế}
\begin{enumerate}
\item \textbf{Bước 1}: Chọn phương trình có hệ số đơn giản ($\pm 1$), biểu diễn một ẩn theo ẩn kia (ví dụ $x = 3 - 2y$).
\item \textbf{Bước 2}: Thế biểu thức vừa tìm được vào phương trình còn lại để được phương trình một ẩn $y$.
\item \textbf{Bước 3}: Giải phương trình một ẩn tìm $y$.
\item \textbf{Bước 4}: Thay giá trị $y$ vừa tìm được vào biểu thức ở Bước 1 để tìm $x$.
\item \textbf{Bước 5}: Kết luận nghiệm của hệ $(x; y)$.
\end{enumerate}

%% DANG 9D1?2-2
\textbf{Dạng 2. Giải hệ bằng phương pháp cộng đại số}
\begin{enumerate}
\item \textbf{Bước 1}: Nhân hai vế của các phương trình với số khác $0$ thích hợp để hệ số của một ẩn trong hai phương trình bằng nhau hoặc đối nhau.
\item \textbf{Bước 2}:
\begin{itemize}
\item Nếu hệ số bằng nhau: lấy từng vế trừ cho nhau.
\item Nếu hệ số đối nhau: lấy từng vế cộng lại với nhau.
\end{itemize}
\item \textbf{Bước 3}: Giải phương trình một ẩn thu được để tìm một biến.
\item \textbf{Bước 4}: Thay giá trị vừa tìm được vào một trong hai phương trình ban đầu để tìm biến còn lại.
\item \textbf{Bước 5}: Kết luận nghiệm $(x; y)$.
\end{enumerate}

%% DANG 9D1?2-3
\textbf{Dạng 3. Hệ phương trình chứa tham số, xác định hệ số}
\begin{enumerate}
\item \textbf{Hệ có nghiệm cho trước}: Thay nghiệm $(x_0; y_0)$ vào hệ để tìm các tham số $m, a, b$.
\item \textbf{Tìm điều kiện tham số để hệ có nghiệm duy nhất}: Giải hệ theo tham số $m$, tìm nghiệm $(x; y)$ theo $m$.
\item Thay $x(m), y(m)$ vào điều kiện đề bài yêu cầu (ví dụ: $x + y > 0$, $x, y$ là các số nguyên\dots) để tìm $m$.
\end{enumerate}

%% DANG 9D1?2-4
\textbf{Dạng 4. Hệ quy về hệ bậc nhất hai ẩn (đặt ẩn phụ)}
\begin{enumerate}
\item \textbf{Bước 1}: Đặt điều kiện xác định cho biểu thức (chú ý ẩn ở mẫu hoặc dưới dấu căn).
\item \textbf{Bước 2}: Nhận diện biểu thức lặp lại, đặt ẩn phụ thích hợp:
$$\text{Ví dụ: } u = \dfrac{1}{x-1},\ v = \dfrac{1}{y+2}\quad \text{hoặc}\quad u = |x|,\ v = \sqrt{y-1}.$$
\item \textbf{Bước 3}: Đưa về hệ phương trình bậc nhất hai ẩn đối với $u, v$ và giải tìm $u, v$.
\item \textbf{Bước 4}: Thay ngược lại tìm $x, y$, đối chiếu điều kiện xác định và kết luận.
\end{enumerate}
